\documentclass[aspectratio=169]{beamer}

% =========================================================
% Кодировка и язык
% =========================================================
\usepackage[utf8]{inputenc}
\usepackage[T2A]{fontenc}
\usepackage[russian]{babel}

% =========================================================
% Математика
% =========================================================
\usepackage{amsmath,amssymb}
\usepackage{bm}

% =========================================================
% Графика и таблицы
% =========================================================
\usepackage{graphicx}
\usepackage{booktabs}
\usepackage{array}

% =========================================================
% TikZ (схемы, архитектуры, диаграммы)
% =========================================================
\usepackage{tikz}
\usetikzlibrary{shapes.geometric, arrows.meta, positioning, calc, fit, backgrounds}

% =========================================================
% Тема
% =========================================================
\usetheme{default}
\usecolortheme{default}

% =========================================================
% Цвета СВФУ
% =========================================================
\definecolor{svfuDark}{RGB}{10,40,90}
\definecolor{svfuAccent}{RGB}{30,144,255}
\definecolor{svfuLight}{RGB}{230,240,255}

% =========================================================
% Стили TikZ (ER-нотация Чена, схемы, уровни, циклы)
% =========================================================
\tikzset{
  entity/.style={
    rectangle, draw=svfuDark, very thick, fill=svfuLight,
    minimum width=2.4cm, minimum height=1cm,
    font=\bfseries\small, align=center, rounded corners=1pt
  },
  relationship/.style={
    diamond, draw=svfuDark, very thick, fill=svfuAccent!25,
    aspect=2.4, font=\scriptsize, align=center, inner sep=1pt
  },
  attribute/.style={
    ellipse, draw=svfuDark, thin, fill=white,
    minimum width=1.8cm, minimum height=0.55cm,
    font=\scriptsize, align=center
  },
  card/.style={font=\tiny, fill=svfuLight, inner sep=1pt},
  layer/.style={
    rectangle, draw=svfuDark, very thick, fill=svfuAccent!15,
    minimum width=6.2cm, minimum height=1.3cm,
    font=\bfseries\normalsize, align=center, rounded corners=2pt
  },
  phase/.style={
    rectangle, draw=svfuDark, thick, fill=svfuLight,
    rounded corners=4pt, minimum width=2.5cm, minimum height=1cm,
    font=\small\bfseries, align=center
  },
  catbox/.style={
    rectangle, draw=svfuDark, thick, fill=white,
    minimum width=2.6cm, minimum height=0.9cm,
    font=\footnotesize, align=center, rounded corners=1pt
  },
  module/.style={
    rectangle, draw=svfuDark, very thick, fill=white,
    minimum width=3.4cm, minimum height=1.1cm,
    font=\footnotesize\bfseries, align=center, rounded corners=2pt
  },
  actor/.style={
    rectangle, draw=svfuDark, thick, fill=svfuLight,
    minimum width=2.6cm, minimum height=1cm,
    font=\small\bfseries, align=center, rounded corners=8pt
  }
}

% =========================================================
% Основные цвета Beamer
% =========================================================
\setbeamercolor{normal text}{fg=black,bg=white}
\setbeamercolor{structure}{fg=svfuDark}
\setbeamercolor{alerted text}{fg=svfuDark}
\setbeamercolor{frametitle}{bg=svfuDark,fg=white}
\setbeamercolor{title}{fg=white}
\setbeamercolor{author}{fg=white}
\setbeamercolor{date}{fg=white}

\setbeamercolor{accent line}{bg=svfuAccent}

\setbeamercolor{footline left}
  {bg=svfuDark,fg=white}

\setbeamercolor{footline center}
  {bg=svfuDark!85,fg=white!80}

\setbeamercolor{footline right}
  {bg=svfuDark!70,fg=white}

\setbeamercolor{block title}
  {bg=svfuDark,fg=white}

\setbeamercolor{block body}
  {bg=svfuLight,fg=black}

% =========================================================
% Шрифты
% =========================================================
\setbeamerfont{frametitle}
  {size=\large,series=\bfseries}

\setbeamerfont{title}
  {size=\Large,series=\bfseries}

% =========================================================
% Убираем навигацию
% =========================================================
\setbeamertemplate{navigation symbols}{}
\setbeamertemplate{headline}{}

% =========================================================
% Заголовок слайда
% =========================================================
\setbeamertemplate{frametitle}{%
  \vspace*{-1pt}%
  \begin{beamercolorbox}[
    wd=\paperwidth,
    ht=3.2ex,
    dp=1.2ex,
    leftskip=10pt,
    sep=0pt
  ]{frametitle}
    \usebeamerfont{frametitle}
    \insertframetitle
  \end{beamercolorbox}%
}

% =========================================================
% Подвал
% =========================================================
\setbeamertemplate{footline}{%
  \begin{beamercolorbox}[
    wd=\paperwidth,
    ht=0.4pt,
    dp=0pt
  ]{accent line}
  \end{beamercolorbox}%

  \hbox to \paperwidth{%

    \begin{beamercolorbox}[
      wd=0.305\paperwidth,
      ht=3ex,
      dp=1.5ex,
      left,
      sep=0pt,
      leftskip=6ex
    ]{footline left}
      \textbf{СВФУ}
    \end{beamercolorbox}%

    \begin{beamercolorbox}[
      wd=0.295\paperwidth,
      ht=3ex,
      dp=1.5ex,
      center,
      sep=0pt
    ]{footline center}
    \end{beamercolorbox}%

    \begin{beamercolorbox}[
      wd=0.4\paperwidth,
      ht=3ex,
      dp=1.5ex,
      right,
      sep=0pt
    ]{footline right}
      \raggedleft
      \today
      \hspace{5ex}
      \insertframenumber\,/\,\inserttotalframenumber
      \hspace{6ex}
    \end{beamercolorbox}%
  }%
}

% =========================================================
% Настройки списков
% =========================================================
\setbeamertemplate{itemize item}
  {\color{svfuAccent}\small$\blacktriangleright$}

\setbeamertemplate{itemize subitem}
  {\color{svfuAccent}\scriptsize$\blacktriangleright$}

\setlength{\leftmargini}{1.5em}

% =========================================================
% Документ
% =========================================================
\begin{document}

% =========================================================
% 1. ТИТУЛ
% =========================================================
\begin{frame}[plain]
  \vspace*{1cm}
  \begin{center}

    \colorbox{svfuDark}{%
      \parbox[c][2.8cm][c]{0.86\paperwidth}{%
        \centering
        {\LARGE\color{white}
          Лекция 2. Системы управления базами данных (СУБД)
        }
      }%
    }

  \end{center}

  \vfill

  \centering
  {\normalsize Дисциплина: Методы и средства проектирования баз данных}\\[4pt]
  {\small
    Составитель: Монастырев Егорий Егорьевич, М-ВТ-25
  }\\[4pt]

  {\small
    Якутск, 2026
  }

  \vfill

\end{frame}


% =========================================================
% СВЯЗЬ С ЛЕКЦИЕЙ 1
% =========================================================
\begin{frame}{Связь с предыдущей лекцией}

На Лекции 1 мы установили:

\begin{itemize}
  \item данные существуют в разных форматах и организуются
        по разным моделям представления (табличной, документной,
        файловой, векторной);
  \item предметная область описывается с помощью ER-модели:
        сущности, атрибуты, связи;
  \item выбор способа хранения — результат анализа данных,
        а не исходная точка проектирования.
\end{itemize}

\vspace{0.4cm}

\begin{alertblock}{Открытый вопрос}
Мы говорили, что данные находятся \textbf{под управлением
специализированного программного обеспечения}.
Что это за программное обеспечение, зачем оно нужно
и как оно устроено?
\end{alertblock}

\end{frame}


% --- Контрольные вопросы лекции ---
\begin{frame}{Ключевые вопросы лекции}

\small
\begin{enumerate}
  \item Чем база данных отличается от СУБД, а СУБД — от системы базы данных?
  \item Какие функции обязана выполнять СУБД и зачем каждая из них нужна?
  \item Что такое трёхуровневая архитектура ANSI/SPARC и что такое независимость данных?
  \item Из каких программных компонентов состоит СУБД изнутри?
  \item Как устроена клиент-серверная модель работы с СУБД?
  \item Какие роли участвуют в работе с СУБД?
  \item Чем настольная СУБД отличается от серверной?
  \item Из каких подъязыков состоит SQL и с какими функциями СУБД они связаны?
  \item Что такое индекс и зачем он нужен, на уровне идеи?
\end{enumerate}

\end{frame}


% =========================================================
% РАЗДЕЛ 1. ОТ БАЗЫ ДАННЫХ К СИСТЕМЕ УПРАВЛЕНИЯ БАЗАМИ ДАННЫХ
% =========================================================

\begin{frame}{Раздел 1. От базы данных к системе управления базами данных}
\centering
\Large
База данных $\;\rightarrow\;$ СУБД $\;\rightarrow\;$ Система базы данных
\end{frame}


\begin{frame}{Почему одной базы данных недостаточно}

Представим, что данные просто лежат в файлах на диске,
без управляющей программы.

\vspace{0.3cm}

\begin{itemize}
  \item Кто гарантирует, что два сотрудника не испортят
        одну и ту же запись, редактируя её одновременно?
  \item Кто восстановит данные после сбоя электропитания
        посреди записи?
  \item Кто проверит, что оценка не выставлена
        несуществующему студенту?
  \item Кто ограничит доступ так, чтобы студент не мог
        править чужие оценки?
\end{itemize}

\vspace{0.3cm}

\begin{alertblock}{Вывод}
Сами данные не решают эти задачи. Ими управляет
специализированное программное обеспечение — \textbf{СУБД}.
\end{alertblock}

\end{frame}


\begin{frame}{База данных: уточнение определения}

\begin{block}{Неполное, бытовое определение}
«База данных — структурированный набор данных.»
\end{block}

\vspace{0.2cm}

\begin{alertblock}{Проблема такого определения}
Под него подходит и обычная таблица Excel или CSV-файл —
но это не база данных в строгом смысле.
\end{alertblock}

\vspace{0.3cm}

\begin{block}{Уточнённое определение}
\textbf{База данных} — организованная в соответствии
с определённой \textbf{схемой} совокупность взаимосвязанных
данных, отражающая состояние некоторой предметной области,
предназначенная для совместного использования многими
приложениями и пользователями и находящаяся под управлением
специализированного программного обеспечения.
\end{block}

\end{frame}


\begin{frame}{СУБД: уточнение определения}

\begin{block}{Неполное, бытовое определение}
«СУБД — программа для создания, хранения и обработки БД.»
\end{block}

\vspace{0.2cm}

\begin{alertblock}{Проблема такого определения}
Такая формулировка не объясняет, \textbf{зачем} именно нужна
СУБД, и создаёт впечатление, что это просто «редактор таблиц».
\end{alertblock}

\vspace{0.3cm}

\begin{block}{Уточнённое определение}
\textbf{СУБД (DBMS)} — комплекс программных средств,
обеспечивающий определение схемы, создание, заполнение,
модификацию и контролируемый доступ к базам данных,
а также реализующий обработку транзакций, поддержание
целостности, разграничение доступа, журналирование
и восстановление данных после сбоев.
\end{block}

\end{frame}


\begin{frame}{От СУБД к системе базы данных}

\begin{center}
\begin{tikzpicture}[node distance=0.32cm and 0.6cm]

\node[entity, minimum width=5.6cm, minimum height=0.75cm] (data) {База данных (данные + схема)};
\node[layer, minimum width=5.6cm, minimum height=0.75cm, below=of data] (dbms) {СУБД};
\node[actor, fill=white, minimum height=0.75cm, below=of dbms, xshift=-1.6cm] (apps) {Приложения};
\node[actor, fill=white, minimum height=0.75cm, below=of dbms, xshift=1.6cm] (people) {Пользователи};

\begin{scope}[on background layer]
\node[draw=svfuAccent, very thick, rounded corners=6pt,
      fit=(data)(dbms)(apps)(people), inner sep=0.3cm,
      label={[font=\bfseries\footnotesize]above:Система базы данных (банк данных)}]
      (frame) {};
\end{scope}

\end{tikzpicture}
\end{center}

\vspace{0.1cm}
\small
\begin{block}{Определение}
\textbf{Система базы данных} — совокупность базы данных,
СУБД, прикладных программ и пользователей, работающих
с этой базой данных как единым целым.
\end{block}
\normalsize

\end{frame}


\begin{frame}{Разграничение понятий}

\small
\begin{table}
\centering
\begin{tabular}{>{\bfseries}p{0.32\textwidth}|p{0.55\textwidth}}
\toprule
Понятие & Что это\\
\midrule
База данных & данные + схема\\
СУБД & программа управления данными\\
Система базы данных & данные + схема + СУБД + приложения + люди\\
Информационная система & система БД + бизнес-логика + интерфейсы + интеграции\\
\bottomrule
\end{tabular}
\end{table}

\vspace{0.4cm}

\begin{alertblock}{На примере «Университета»}
Данные о студентах и группах — \textbf{база данных};
PostgreSQL, под управлением которого они хранятся, —
\textbf{СУБД}; портал деканата, объединяющий БД, СУБД,
приложение и сотрудников, — \textbf{система базы данных}.
\end{alertblock}

\end{frame}


% =========================================================
% РАЗДЕЛ 2. НАЗНАЧЕНИЕ И ФУНКЦИИ СУБД
% =========================================================

\begin{frame}{Раздел 2. Назначение и функции СУБД}
\centering
\Large
Зачем нужна СУБД и что именно она делает?
\end{frame}


\begin{frame}{Общее назначение СУБД}

СУБД решает задачи, которые невозможно надёжно решить
средствами отдельных файлов и прикладных программ:

\vspace{0.3cm}

\begin{itemize}
  \item предоставляет единый, контролируемый способ
        доступа к данным;
  \item берёт на себя ответственность за корректность
        данных при сбоях и параллельной работе;
  \item скрывает от приложений детали физического
        хранения данных.
\end{itemize}

\vspace{0.3cm}

\begin{block}{Идея}
Приложения работают с данными \textbf{через} СУБД,
а не напрямую с файлами на диске.
\end{block}

\end{frame}


\begin{frame}{Функции СУБД}

\small
\begin{table}
\centering
\begin{tabular}{p{0.36\textwidth}|p{0.5\textwidth}}
\toprule
\textbf{Функция} & \textbf{Суть}\\
\midrule
Хранение данных & организация данных на диске в виде таблиц, файлов, страниц\\
Обработка и оптимизация запросов & разбор и эффективное выполнение запросов к данным\\
Управление транзакциями & выполнение операций по принципу «всё или ничего»\\
Контроль целостности & проверка правил, которым должны удовлетворять данные\\
Управление многопользовательским доступом & корректная одновременная работа многих пользователей\\
Журналирование и восстановление & защита от потери данных при сбоях\\
Безопасность & разграничение прав доступа к данным\\
\bottomrule
\end{tabular}
\end{table}

\end{frame}


\begin{frame}{Функции СУБД: пример на «Университете»}

\small
\begin{itemize}
  \item \textbf{Целостность} — нельзя выставить оценку студенту
        с несуществующим номером;
  \item \textbf{Многопользовательский доступ} — два сотрудника
        деканата одновременно редактируют оценки одного
        студента, и это не приводит к потере данных;
  \item \textbf{Журналирование и восстановление} — при сбое
        питания посреди сохранения ведомости данные
        не теряются;
  \item \textbf{Безопасность} — студент видит только свои
        оценки, а не оценки всей группы;
  \item \textbf{Транзакции} — перевод студента из одной группы
        в другую либо выполняется полностью, либо не
        выполняется вовсе.
\end{itemize}

\end{frame}


\begin{frame}{От функций СУБД — к свойствам базы данных}

\begin{center}
\begin{tikzpicture}[node distance=0.22cm and 1.2cm]

\node[module, fill=svfuLight, minimum height=0.68cm] (f1) {Целостность};
\node[module, fill=svfuLight, minimum height=0.68cm, below=of f1] (f2) {Многопольз.\\доступ};
\node[module, fill=svfuLight, minimum height=0.68cm, below=of f2] (f3) {Безопасность};
\node[module, fill=svfuLight, minimum height=0.68cm, below=of f3] (f4) {Хранение};

\node[catbox, minimum height=0.68cm, right=2.4cm of f1] (p1) {Целостность (БД)};
\node[catbox, minimum height=0.68cm, right=2.4cm of f2] (p2) {Разделяемость};
\node[catbox, minimum height=0.68cm, right=2.4cm of f3] (p3) {Защищённость};
\node[catbox, minimum height=0.68cm, right=2.4cm of f4] (p4) {Интегрированность};

\draw[-{Latex[length=2mm]}, thick, svfuAccent] (f1) -- (p1);
\draw[-{Latex[length=2mm]}, thick, svfuAccent] (f2) -- (p2);
\draw[-{Latex[length=2mm]}, thick, svfuAccent] (f3) -- (p3);
\draw[-{Latex[length=2mm]}, thick, svfuAccent] (f4) -- (p4);

\node[above=0.15cm of f1, font=\footnotesize\bfseries] {Функция СУБД};
\node[above=0.15cm of p1, font=\footnotesize\bfseries] {Свойство БД};

\end{tikzpicture}
\end{center}

\small
\begin{alertblock}{Главная мысль}
Свойства базы данных, названные в Лекции 1
(целостность, разделяемость, защищённость,
интегрированность), — не сами по себе, а \textbf{результат}
работы конкретных функций СУБД.
\end{alertblock}
\normalsize

\end{frame}


% =========================================================
% РАЗДЕЛ 3. АРХИТЕКТУРА ПРЕДСТАВЛЕНИЯ ДАННЫХ
% =========================================================

\begin{frame}{Раздел 3. Архитектура представления данных: ANSI/SPARC}
\centering
\Large
Как одни и те же данные видны на разных уровнях абстракции?
\end{frame}


\begin{frame}{Схема и экземпляр базы данных}

\begin{itemize}
  \item \textbf{Схема базы данных} — формальное описание
        структуры БД: таблицы, атрибуты, типы данных,
        связи, ограничения. Меняется редко.
  \item \textbf{Экземпляр базы данных} — конкретное состояние
        данных в определённый момент времени. Меняется
        практически при каждой операции записи.
\end{itemize}

\vspace{0.4cm}

\begin{alertblock}{Аналогия для запоминания}
Схема — это «тип» (структура), экземпляр — это «значение»
(конкретное содержимое) в данный момент.
\end{alertblock}

\vspace{0.2cm}
\footnotesize
Пример: схема таблицы \texttt{student} (id, ФИО, group\_id)
не меняется годами; строки этой таблицы меняются
каждый день.

\end{frame}


\begin{frame}{Трёхуровневая архитектура ANSI/SPARC}

Это не этапы проектирования, а архитектурная модель
самой СУБД: способ, которым данные одновременно
представлены на трёх уровнях.

\begin{center}
\begin{tikzpicture}[node distance=0.6cm]
\node[layer] (ext) {Внешний уровень\\ \footnotesize представления пользователей};
\node[layer, below=of ext] (con) {Концептуальный уровень\\ \footnotesize логическая структура БД};
\node[layer, below=of con] (int) {Внутренний уровень\\ \footnotesize физическое хранение};
\draw[-{Latex[length=2.5mm]}, thick, svfuAccent] (ext) -- (con);
\draw[-{Latex[length=2.5mm]}, thick, svfuAccent] (con) -- (int);
\end{tikzpicture}
\end{center}

\begin{center}
Цель — \textbf{независимость данных}.
\end{center}

\end{frame}


\begin{frame}{Три уровня подробнее}

\begin{itemize}
  \item \textbf{Внешний} — как данные видит конкретный
        пользователь или приложение (студент видит свои
        оценки, деканат — статистику факультета).
  \item \textbf{Концептуальный} — единое логическое описание
        всей БД: сущности, атрибуты, связи, ограничения.
  \item \textbf{Внутренний} — как данные физически хранятся:
        файлы, страницы, индексы, буферы.
\end{itemize}

\vspace{0.3cm}

\begin{alertblock}{Уточнение}
Соответствие «внешний уровень $\approx$ \texttt{VIEW}
в PostgreSQL» — лишь приближённое: полноценный внешний
уровень формируется также правами доступа (\texttt{GRANT}/
\texttt{REVOKE}) и структурой запросов приложения,
а не только объектами \texttt{VIEW}.
\end{alertblock}

\end{frame}


\begin{frame}{Независимость данных}

\begin{table}
\centering
\begin{tabular}{p{0.3\textwidth}|p{0.58\textwidth}}
\toprule
\textbf{Тип} & \textbf{Что позволяет}\\
\midrule
Физическая &
изменять способы хранения и индексы
без переписывания приложений\\
Логическая &
изменять концептуальную схему,
сохраняя работоспособность существующих представлений\\
\bottomrule
\end{tabular}
\end{table}

\vspace{0.4cm}

\begin{alertblock}{Важная оговорка}
Логическая независимость имеет пределы: удаление или
переименование атрибута, который уже использовало
приложение, всё равно нарушит его работу.
\end{alertblock}

\end{frame}


% =========================================================
% РАЗДЕЛ 4. ВНУТРЕННЯЯ АРХИТЕКТУРА СУБД
% =========================================================

\begin{frame}{Раздел 4. Внутренняя архитектура СУБД: программные компоненты}
\centering
\Large
Из чего состоит СУБД изнутри?
\end{frame}


\begin{frame}{ANSI/SPARC и внутренняя архитектура — не одно и то же}

\begin{block}{Важное разграничение}
Архитектура ANSI/SPARC отвечает на вопрос
\textbf{«как данные видны снаружи»}.\\[4pt]
Внутренняя архитектура СУБД отвечает на вопрос
\textbf{«из каких программных модулей СУБД состоит изнутри»}.
\end{block}

\vspace{0.4cm}

Это разные по природе понятия, которые не стоит путать:
первое — про уровни абстракции представления данных,
второе — про устройство программного ядра.

\end{frame}


\begin{frame}{Компоненты ядра СУБД}

\small
\begin{table}
\centering
\begin{tabular}{p{0.42\textwidth}|p{0.44\textwidth}}
\toprule
\textbf{Компонент} & \textbf{Что делает}\\
\midrule
Обработчик запросов & разбор и оптимизация запросов\\
Менеджер транзакций & контроль атомарности и изоляции операций\\
Менеджер хранения и буфера & управление файлами, страницами, кэшем в ОЗУ\\
Менеджер журнала и восстановления & фиксация изменений, восстановление после сбоев\\
Словарь данных / системный каталог & хранение метаданных о структуре БД\\
\bottomrule
\end{tabular}
\end{table}

\vspace{0.3cm}
\footnotesize
Каждый компонент отвечает за одну из функций СУБД,
рассмотренных в разделе 2.

\end{frame}


\begin{frame}{Путь одного запроса через компоненты СУБД}

\begin{center}
\resizebox{0.98\textwidth}{!}{%
\begin{tikzpicture}[node distance=0.55cm and 0.7cm]

\node[actor, fill=svfuLight] (app) {Приложение\\\texttt{SELECT ...}};
\node[module, right=of app] (parser) {Обработчик\\запросов};
\node[module, right=of parser] (tx) {Менеджер\\транзакций};
\node[module, right=of tx] (storage) {Менеджер\\хранения};
\node[entity, right=of storage, minimum height=1.1cm] (disk) {Данные\\на диске};

\draw[-{Latex[length=2mm]}, thick, svfuAccent] (app) -- (parser);
\draw[-{Latex[length=2mm]}, thick, svfuAccent] (parser) -- (tx);
\draw[-{Latex[length=2mm]}, thick, svfuAccent] (tx) -- (storage);
\draw[-{Latex[length=2mm]}, thick, svfuAccent] (storage) -- (disk);

\node[module, below=1.1cm of tx] (log) {Менеджер\\журнала};
\draw[-{Latex[length=2mm]}, thick, svfuDark, dashed] (tx) -- (log);

\node[catbox, below=1.1cm of parser] (dict) {Словарь\\данных};
\draw[-{Latex[length=2mm]}, thick, svfuDark, dashed] (parser) -- (dict);

\end{tikzpicture}%
}
\end{center}

\footnotesize
Схема иллюстрирует общий путь запроса без разбора
алгоритмов оптимизации: обработчик обращается к словарю
данных, чтобы узнать структуру таблиц; менеджер транзакций
взаимодействует с менеджером журнала для гарантии
устойчивости изменений.

\end{frame}


% =========================================================
% РАЗДЕЛ 5. МНОГОПОЛЬЗОВАТЕЛЬСКИЙ РЕЖИМ
% =========================================================

\begin{frame}{Раздел 5. Многопользовательский режим и участники работы с СУБД}
\centering
\Large
Как СУБД обслуживает многих пользователей одновременно?
\end{frame}


\begin{frame}{Клиент-серверная модель доступа к СУБД}

\begin{center}
\begin{tikzpicture}[node distance=0.4cm and 0.9cm]

\node[actor, fill=white, minimum width=2.4cm] (c1) {Деканат};
\node[actor, fill=white, minimum width=2.4cm, below=0.35cm of c1] (c2) {Приёмная комиссия};
\node[actor, fill=white, minimum width=2.4cm, below=0.35cm of c2] (c3) {Личный кабинет};

\node[layer, right=1.8cm of c2, minimum width=3.6cm] (server) {Сервер СУБД\\\footnotesize (PostgreSQL)};

\node[entity, right=1.8cm of server, minimum height=1.4cm, minimum width=2.2cm] (files) {Файлы\\базы данных};

\draw[thick, svfuAccent] (c1) -- (server);
\draw[thick, svfuAccent] (c2) -- (server);
\draw[thick, svfuAccent] (c3) -- (server);
\draw[-{Latex[length=2.5mm]}, thick, svfuDark] (server) -- (files);

\end{tikzpicture}
\end{center}

\begin{alertblock}{Идея}
Множество клиентских приложений одновременно подключаются
к \textbf{одному} серверу СУБД, а не к файлам напрямую.
\end{alertblock}

\end{frame}


\begin{frame}{Параллельная работа: постановка проблемы}

\begin{center}
\large
Сотрудник А редактирует оценку студента 101
\end{center}
\begin{center}
\large
Сотрудник Б \textbf{одновременно} редактирует ту же оценку
\end{center}

\vspace{0.4cm}

\begin{itemize}
  \item Что произойдёт, если оба сохранят изменения
        одновременно?
  \item Как СУБД не допускает потери одного из изменений?
\end{itemize}

\vspace{0.3cm}

\begin{alertblock}{На этом этапе — только постановка вопроса}
Сама механика решения (блокировки, уровни изоляции)
рассматривается в отдельной будущей лекции
о транзакциях и параллельном выполнении.
Сейчас важно зафиксировать: СУБД обязана
это контролировать.
\end{alertblock}

\end{frame}


\begin{frame}{Роли участников работы с СУБД}

\small
\begin{table}
\centering
\begin{tabular}{>{\bfseries}p{0.32\textwidth}|p{0.55\textwidth}}
\toprule
Роль & Что делает\\
\midrule
Администратор БД (АБД) &
администрирует безопасность, резервное копирование
и восстановление, следит за производительностью\\
Прикладные программисты &
используют транзакции и целостность через приложение,
не работая с ядром СУБД напрямую\\
Конечные пользователи &
взаимодействуют с данными через приложение,
не видя внутреннего устройства СУБД\\
\bottomrule
\end{tabular}
\end{table}

\end{frame}


% =========================================================
% РАЗДЕЛ 6. КЛАССИФИКАЦИЯ СУБД
% =========================================================

\begin{frame}{Раздел 6. Классификация СУБД и средства работы с ней}
\centering
\Large
Какие бывают СУБД и как к ним подключаются?
\end{frame}


\begin{frame}{Классификация СУБД по модели данных (напоминание)}

Из Лекции 1: способ представления данных определяет
и тип СУБД, работающей с этими данными.

\vspace{0.4cm}

\begin{table}
\centering
\small
\begin{tabular}{p{0.3\textwidth}|p{0.3\textwidth}|p{0.3\textwidth}}
\toprule
Реляционные & Документные & Другие\\
\midrule
PostgreSQL, MySQL, Oracle & MongoDB & Vector DB и т.д.\\
\bottomrule
\end{tabular}
\end{table}

\vspace{0.4cm}

\begin{alertblock}{Далее в этой лекции}
Речь пойдёт о другой оси классификации — не о модели
данных, а об архитектуре развёртывания СУБД.
\end{alertblock}

\end{frame}


\begin{frame}{Настольные и серверные СУБД}

\small
\begin{table}
\centering
\begin{tabular}{p{0.28\textwidth}|p{0.28\textwidth}|p{0.28\textwidth}}
\toprule
 & \textbf{Access} & \textbf{PostgreSQL}\\
\midrule
Тип & настольная СУБД & серверная СУБД\\
Интерфейс & сильный GUI & SQL + инструменты\\
Масштаб & небольшие системы & крупные системы\\
Многопольз. доступ & ограничен & полноценный\\
\bottomrule
\end{tabular}
\end{table}

\normalsize

\vspace{0.4cm}

\begin{block}{Вывод}
Access удобен для локальных обучающих задач;
PostgreSQL — основная серверная СУБД курса.
\end{block}

\end{frame}


\begin{frame}{Инструменты администрирования серверной СУБД}

Для подключения к серверной СУБД используются
отдельные клиентские инструменты:

\vspace{0.4cm}

\begin{itemize}
  \item \textbf{pgAdmin} — графическое администрирование
        PostgreSQL: просмотр таблиц, выполнение запросов,
        настройка прав доступа;
  \item \textbf{psql} — консольный (CLI) клиент PostgreSQL
        для прямого взаимодействия с сервером.
\end{itemize}

\vspace{0.4cm}

\begin{alertblock}{Важно}
Это средства \textbf{работы} с уже существующей серверной
СУБД, а не средства проектирования схемы БД
(те рассматривались в Лекции 1: draw.io, pgModeler, DBeaver).
\end{alertblock}

\end{frame}


% =========================================================
% РАЗДЕЛ 7. SQL КАК СРЕДСТВО РАБОТЫ С СУБД
% =========================================================

\begin{frame}{Раздел 7. SQL как средство работы с СУБД (обзорно)}
\centering
\Large
Как практически формулируется команда для работы с СУБД?
\end{frame}


\begin{frame}{SQL как декларативный язык}

\begin{block}{Определение}
\textbf{SQL} (Structured Query Language) — стандартизированный
декларативный язык для определения структуры данных,
выполнения операций над данными и управления доступом
в реляционных СУБД.
\end{block}

\vspace{0.3cm}

\begin{itemize}
  \item декларативный — пользователь описывает,
        \textbf{что} нужно получить, а не \textbf{как} именно
        это будет вычислено;
  \item поддерживается практически всеми реляционными СУБД
        (с диалектными отличиями).
\end{itemize}

\vspace{0.3cm}

\begin{alertblock}{На этом этапе курса}
Подробный синтаксис запросов не рассматривается —
только структура языка и её связь с функциями СУБД.
\end{alertblock}

\end{frame}


\begin{frame}{Подъязыки SQL и функции СУБД}

\small
\begin{table}
\centering
\begin{tabular}{>{\bfseries}p{0.14\textwidth}|p{0.36\textwidth}|p{0.36\textwidth}}
\toprule
Подъязык & Назначение & Пример команд\\
\midrule
DDL & определение структуры данных & \texttt{CREATE}, \texttt{ALTER}, \texttt{DROP}\\
DML & работа с данными & \texttt{SELECT}, \texttt{INSERT}, \texttt{UPDATE}, \texttt{DELETE}\\
DCL & управление доступом & \texttt{GRANT}, \texttt{REVOKE}\\
TCL & управление транзакциями & \texttt{COMMIT}, \texttt{ROLLBACK}\\
\bottomrule
\end{tabular}
\end{table}

\vspace{0.4cm}

\begin{alertblock}{Главная идея}
SQL — не хаотичный набор команд, а язык, структурно
повторяющий функции самой СУБД: хранение и структуру
(DDL), обработку данных (DML), безопасность (DCL),
транзакции (TCL).
\end{alertblock}

\end{frame}


% =========================================================
% РАЗДЕЛ 8. ФИЗИЧЕСКОЕ ХРАНЕНИЕ ДАННЫХ
% =========================================================

\begin{frame}{Раздел 8. Физическое хранение данных (обзорно)}
\centering
\Large
Где физически лежат данные, с которыми работает SQL?
\end{frame}


\begin{frame}{От логической модели к физическому хранению}

\begin{center}
\resizebox{0.98\textwidth}{!}{%
\begin{tikzpicture}[node distance=1cm]

\node[entity, minimum width=3.1cm] (log) {Логическая структура\\(таблицы, столбцы)};
\node[catbox, right=of log] (files) {Файлы};
\node[catbox, right=of files] (pages) {Страницы};
\node[catbox, right=of pages, minimum width=2.6cm] (buf) {Буфер\\(кэш в ОЗУ)};

\draw[-{Latex[length=2.5mm]}, thick, svfuAccent] (log) -- (files);
\draw[-{Latex[length=2.5mm]}, thick, svfuAccent] (files) -- (pages);
\draw[-{Latex[length=2.5mm]}, thick, svfuAccent] (pages) -- (buf);

\end{tikzpicture}%
}
\end{center}

\footnotesize
Пользователь и приложение работают только с логической
структурой; всё, что ниже, скрыто внутренним уровнем
архитектуры ANSI/SPARC (раздел 3).

\end{frame}


\begin{frame}{Файлы, страницы, буферы — кратко}

\begin{itemize}
  \item \textbf{Файл} — единица хранения таблицы или её части
        на диске;
  \item \textbf{Страница} — блок фиксированного размера,
        которым СУБД читает и записывает данные на диск;
  \item \textbf{Буфер} — область оперативной памяти,
        где хранятся недавно использованные страницы,
        чтобы не обращаться к диску повторно.
\end{itemize}

\vspace{0.3cm}

\begin{alertblock}{Важно}
Это обзорное представление. Алгоритмы работы буфера
и физической организации файлов — предмет отдельной
будущей лекции о физическом проектировании.
\end{alertblock}

\end{frame}


\begin{frame}{Индекс как идея}

Поиск студента по фамилии в таблице из миллиона строк:

\vspace{0.3cm}

\begin{itemize}
  \item \textbf{без индекса} — СУБД просматривает
        все строки подряд;
  \item \textbf{с индексом} — СУБД использует
        вспомогательную структуру, которая сразу указывает,
        где искать нужные строки.
\end{itemize}

\vspace{0.4cm}

\begin{block}{Определение}
\textbf{Индекс} — вспомогательная структура данных,
ускоряющая поиск записей ценой дополнительного
объёма хранения и дополнительных затрат при изменении
данных.
\end{block}

\vspace{0.2cm}
\footnotesize
Синтаксис создания индекса и стратегии индексирования —
предмет отдельной будущей лекции.

\end{frame}


% =========================================================
% РАЗДЕЛ 9. ИТОГИ
% =========================================================

\begin{frame}{Раздел 9. Итоги лекции и переход к следующей теме}
\centering
\Large
От файла — к системе управления базами данных
\end{frame}


\begin{frame}{Итоговая схема: от файла к системе базы данных}

\begin{center}
\resizebox{0.98\textwidth}{!}{%
\begin{tikzpicture}[node distance=0.5cm and 0.4cm]

\node[phase] (file) {Файл};
\node[phase, right=of file] (table) {Таблица};
\node[phase, right=of table] (db) {База\\данных};
\node[phase, right=of db] (dbms) {СУБД};
\node[phase, right=of dbms] (sys) {Система\\базы данных};

\draw[-{Latex[length=2.5mm]}, thick, svfuAccent] (file) -- (table);
\draw[-{Latex[length=2.5mm]}, thick, svfuAccent] (table) -- (db);
\draw[-{Latex[length=2.5mm]}, thick, svfuAccent] (db) -- (dbms);
\draw[-{Latex[length=2.5mm]}, thick, svfuAccent] (dbms) -- (sys);

\end{tikzpicture}%
}
\end{center}

\vspace{0.4cm}

\begin{alertblock}{Главная идея лекции}
СУБД — не просто программа для хранения таблиц,
а комплекс функций (транзакции, целостность,
многопользовательский доступ, восстановление,
безопасность), реализованных набором программных
компонентов и доступных через SQL.
\end{alertblock}

\end{frame}


\begin{frame}{Что будет раскрыто подробнее в следующих лекциях}

\small
\begin{table}
\centering
\begin{tabular}{p{0.5\textwidth}|p{0.4\textwidth}}
\toprule
\textbf{Тема} & \textbf{Где}\\
\midrule
Реляционная модель, ключи, ограничения целостности (SQL) & Лекция 3\\
Нормализация и нормальные формы & Лекция 3 / отдельная лекция\\
Полный синтаксис SQL, \texttt{JOIN}, агрегатные функции & SQL-практикум\\
Механика транзакций: изоляция, блокировки & Отдельная лекция о транзакциях\\
Детали журналирования и восстановления & Та же лекция о надёжности\\
Синтаксис и стратегии индексирования & Лекция о физическом проектировании\\
\bottomrule
\end{tabular}
\end{table}

\end{frame}


\begin{frame}{Контрольные вопросы}

\small
\begin{enumerate}
  \item В чём разница между базой данных, СУБД и системой базы данных?
  \item Почему определение «СУБД — программа для работы с БД» неполно?
  \item Что различают схема и экземпляр базы данных?
  \item Из каких трёх уровней состоит архитектура ANSI/SPARC?
  \item Чем архитектура ANSI/SPARC отличается от внутренней архитектуры СУБД?
  \item Назовите пять программных компонентов ядра СУБД.
  \item Что такое клиент-серверная модель работы с СУБД?
  \item Чем отличаются DDL, DML, DCL и TCL?
  \item Зачем нужен индекс и какова его цена?
\end{enumerate}

\end{frame}


\begin{frame}{Задание к следующей лекции}

Возьмите предметную область и ER-диаграмму,
разработанную к Лекции 1.

\vspace{0.4cm}

Определите:
\begin{enumerate}
  \item какие функции СУБД будут задействованы при работе
        с вашими данными (целостность, многопользовательский
        доступ и т.\,д.) — приведите по одному примеру
        на функцию;
  \item какие сущности вашей предметной области логично
        превратить в таблицы реляционной СУБД.
\end{enumerate}

\vspace{0.4cm}

\textbf{Результат:} короткий список функций СУБД
с примерами + перечень будущих таблиц.

\end{frame}


% --- Спасибо ---
\begin{frame}[plain,noframenumbering]

\thispagestyle{empty}

\vspace*{\fill}

\begin{center}

{\LARGE\bfseries
Спасибо за внимание!
}

\end{center}

\vspace*{\fill}

\end{frame}


\end{document}
